% related-work.tex -- BNR-2026-013

\section{Related Work}
\label{sec:related}

\subsection{ZK Proof Systems}

Groth16~\cite{groth16} remains the most widely deployed SNARK for
on-chain verification due to its minimal proof size (128 bytes
compressed on BN254) and low verification cost ($\sim$220K
gas~\cite{nebra-gas}). However, the per-circuit trusted setup---a
two-phase ceremony where Phase~1 is universal but Phase~2 is
circuit-specific---creates operational friction for evolving systems.

PLONK~\cite{plonk} introduced universal and updateable SRS via
polynomial commitment schemes, enabling a single Powers-of-Tau
ceremony to support arbitrary circuits up to a maximum size. The
PLONKish arithmetization generalizes R1CS by allowing custom
polynomial gates of arbitrary degree, copy constraints via
permutation arguments, and lookup tables via
plookup~\cite{plookup,plonkup}.

The halo2 proving system implements PLONKish arithmetization with
configurable polynomial commitment backends. The original Zcash
implementation uses Inner Product Argument (IPA) on Pasta
curves~\cite{halo}, achieving transparent setup but at the cost of
EVM incompatibility (no precompile for Pallas/Vesta curves).
The Privacy and Scaling Explorations (PSE) fork~\cite{pse-halo2}
adapted halo2 to KZG commitments on BN254, enabling efficient EVM
verification via existing pairing precompiles.

FFLONK~\cite{fflonk} optimizes the PLONK verifier using
Fast-Fourier techniques, achieving $\sim$200K gas verification
cost~\cite{orbiter-gas}---comparable to Groth16---at the expense
of implementation complexity.

Plonky2~\cite{plonky2} combines PLONK with FRI (Fast Reed-Solomon
Interactive Oracle Proofs) over the Goldilocks field, achieving
170~ms recursive proofs but producing 43--130~KB proofs unsuitable
for direct on-chain verification.

Nova~\cite{nova} and SuperNova~\cite{supernova} introduce
Incrementally Verifiable Computation (IVC) via folding schemes,
offering an alternative recursion model with $\sim$10K constraint
verifier circuits.

\subsection{ZK-Friendly Hash Functions}

Poseidon~\cite{poseidon} is a purpose-designed hash function for
ZK circuits, requiring $\sim$300 R1CS constraints per invocation
compared to 25,000--30,000 for SHA-256. Poseidon2~\cite{poseidon2}
further reduces round counts. The Basis Network state trie uses
Poseidon over BN254, making custom gate optimization for this hash
function particularly impactful.

\subsection{Production zkEVM Deployments}

Several production systems inform our analysis:

\textbf{Scroll}~\cite{scroll} pioneered halo2-KZG for zkEVM,
employing 2000+ custom gates for EVM opcode verification on BN254.
Scroll has since begun migrating to OpenVM for improved modularity.

\textbf{Axiom}~\cite{axiom} deploys a production halo2-KZG prover
on Ethereum mainnet with a fixed 420K gas verification cost,
demonstrating the viability of this approach for high-value
applications.

\textbf{Polygon zkEVM}~\cite{polygon-zkevm} uses a two-layer
architecture: eSTARK for fast inner proving with FFLONK for
compact L1 verification at $\sim$179K gas.

\textbf{zkSync Era}~\cite{zksync} employs a custom PLONK variant
(Boojum) with FRI over the Goldilocks field, wrapping into
PLONK-KZG on BN254 for L1 verification.

\textbf{Taiko}~\cite{taiko} uses PSE halo2-KZG directly for
Type~1 zkEVM proving, validating the approach for full EVM
equivalence.

\subsection{Benchmark Studies}

The zk-Bench framework~\cite{zkbench} provides cross-system
benchmarking methodology. Analysis of zkEVM constraint-level
architectures~\cite{zkevm-constraints} documents custom gate
design patterns across production systems. Research on ZK rollup
performance bottlenecks~\cite{zk-bottlenecks} identifies proving
time and circuit complexity as dominant cost factors. The
Plonkify transpiler~\cite{plonkify} demonstrates 2.35$\times$
constraint inflation when converting R1CS to PLONKish without
custom gates, motivating the use of native PLONKish circuit
design.
